\begin{enumerate}
    %%%---2a---%%%
    \item Yes, $x_1^* = 0$, $x_2^* = 3$ is a feasible solution. The verification is done by substitution:
    \begin{align}
     0 + 3 & = 3 \qquad \checkmark \nonumber  \\
     0 & \geq 0 \qquad \checkmark \nonumber  \\
     3 & \geq 0. \qquad \checkmark \nonumber 
    \end{align}
%
    The first two constraints are active at $\{x_1^*, x_2^*\}$. In
    particular, notice that equality constraints are always active
    when verified. The  point is labeled $A$ in the figure below.

    \item No, $x_1^* = 2$, $x_2^* = 2$ is not a feasible solution, since the equality constraint is not verified:
    \begin{align}
     2 + 2 & = 3 \qquad \xmark \nonumber  \\
     2 & \geq 0 \qquad \checkmark \nonumber \\
     2 & \geq 0. \qquad \checkmark \nonumber
    \end{align}
    
    The point is labeled $B$ in the figure below.
    
\end{enumerate}

The feasible region for the problem is the segment represented in bold in the following graph.
     
     \begin{center}
     \begin{tikzpicture}[domain=-0.6:4.6,range=-0.5:3.6]
      
      \draw[->] (-0.2,0) -- (4.2,0) node[right] {$x_1$};
      \draw[->] (0,-0.2) -- (0,3.7) node[above] {$x_2$};
      
      \draw[color=black] (-0.6,3.6) -- (3.4,-0.4);
      \draw[color=black, ultra thick] (0,3) -- (3,0);
      \node at (3.6,-0.6) {$x_1 + x_2 = 3$};
      
      \node [circle,fill=black,draw,scale=0.2] at (0,3) {A};
      \node at (-0.2,2.8) {A};
      \node [circle,fill=black,draw,scale=0.2] at (2,2) {B};
      \node at (2.2,2.2) {B};
    
    \end{tikzpicture}
    \end{center}
